%% OEIS Program Synthesis, Verified Semantics, and Graded Similarity
%%
%% A collecting account of the OEIS-synthesis line of work: a verified domain-specific
%% language for integer-sequence programs, a neural synthesizer whose contracts are the
%% prover's theorems, and a graded (probabilistic-logic) similarity layer over the corpus.
%%
%% Build:
%%   pdflatex OEIS.tex

\documentclass[11pt]{article}
\usepackage[T1]{fontenc}
\usepackage[utf8]{inputenc}
\usepackage{lmodern}
\usepackage{geometry}
\geometry{margin=1in}
\usepackage{hyperref}
\usepackage{booktabs}
\usepackage{longtable}
\usepackage{listings}
\usepackage{array}
\usepackage{amsmath}
\usepackage{amssymb}
\usepackage{amsthm}
\setlength{\emergencystretch}{3em}

\lstset{
  basicstyle=\ttfamily\small,
  breaklines=true,
  frame=single,
  columns=fullflexible
}

\newcommand{\code}[1]{\texttt{#1}}
\newcommand{\sig}{\ensuremath{\sigma}}
\newtheorem{obligation}{Obligation}
\newtheorem{property}{Property}

\title{OEIS Program Synthesis, Verified Semantics, and Graded Similarity:\\
       A Unified Account}
\author{Zarathustra J.\ Goertzel}
\date{\today}

\begin{document}
\maketitle

\begin{abstract}
We collect an ongoing line of work built around the On-Line Encyclopedia of Integer
Sequences (OEIS) as a testbed for verified program synthesis. Three strands meet here.
First, a small domain-specific language (DSL) of integer-sequence programs is given a
machine-checked operational semantics in Lean, together with a token \emph{recognizer},
a big-step generative-semantics (GSLT) layer with modal emission specifications, and a
family of formal \emph{obligations} that a synthesizer must respect. Second, a neural
sequence-to-graph synthesizer is organized so that the network's operational
contracts---a legality mask, a canonical form, relabeling invariance---are exactly the
theorems the prover discharges; the architecture is read as a specification. Third, a
graded similarity layer expressed in probabilistic-logic truth values $\langle
\text{strength},\text{confidence}\rangle$ recovers the structure that an exact
deduplication hash discards: it audits train/test leakage and, via a \emph{sound abstract
interpreter}, attaches provable properties (totality, sign, parity) to each program. We
report concrete measurements over a corpus of $104{,}083$ synthesized programs, and we
argue that the same provable-property machinery serves two roles in an OEIS synthesis
loop: it \emph{filters} candidate programs and it \emph{measures} them theoretically,
turning a leakage audit into a verified conjecture generator.
\end{abstract}

\section{Introduction}

The OEIS is a natural benchmark for program synthesis: each entry is an integer sequence,
and a solution is a program that reproduces its terms. Recent self-learning systems
synthesize such programs at scale from a small combinator DSL, discovering programs for a
large fraction of the OEIS without human-written training
data~\cite{gauthier2023,gauthier2025}. Our contribution is not another synthesizer but a
\emph{verified and measured} account of the objects a synthesizer manipulates.

We separate three concerns that are often conflated:
\begin{itemize}
  \item \textbf{What a program \emph{is}} --- its syntax over a fixed operator vocabulary,
        and the machine-checked semantics of running it (Sections~\ref{sec:dsl}--\ref{sec:gslt}).
  \item \textbf{What a program \emph{does}} --- its input/output behaviour on a finite probe
        window (its \emph{signature} \sig), the extensional object
        (Section~\ref{sec:extensional}).
  \item \textbf{What a program \emph{provably has} as properties} --- totality, sign,
        parity, growth: the intensional object, obtained by sound static analysis rather
        than read off finitely many outputs (Section~\ref{sec:intensional}).
\end{itemize}
The payoff of keeping these apart is a clean division of labour: the extensional layer
gives cheap graded evidence but can never certify behaviour beyond the probe window; the
intensional layer certifies properties for all inputs but is conservative. Combined, they
sharpen both benchmark hygiene and the discovery of conjectures.

\section{The OEIS synthesis DSL}
\label{sec:dsl}

Programs are postfix token streams over a $16$-operator vocabulary evaluated on a
two-register integer state. The operators are constants \code{zero}, \code{one},
\code{two}; the index variable \code{x} $=n$ and an inner variable \code{y}; binary
arithmetic \code{addi}, \code{diff}, \code{mult}, \code{divi}, \code{modu}; a conditional
\code{cond} (test $\le 0$); bounded iterators \code{loop} and \code{loop2}; an unbounded
search \code{compr} (the $a$-th $n\ge 0$ satisfying a predicate); and stack operators
\code{push}, \code{pop}. Table~\ref{tab:ops} summarizes the arities and the partiality
character of each operator, which is decisive for the analysis of
Section~\ref{sec:intensional}.

\begin{table}[t]
\centering
\begin{tabular}{@{}llll@{}}
\toprule
operator & arity & role & partiality / termination \\
\midrule
\code{zero},\code{one},\code{two} & $0$ & constants & total \\
\code{x},\code{y} & $0$ & variables ($n\ge 0$; inner $\ge 0$) & total \\
\code{addi},\code{diff},\code{mult} & $2$ & arithmetic & total (resource: overflow) \\
\code{divi},\code{modu} & $2$ & division / modulo & \emph{genuine}: error on zero divisor \\
\code{cond} & $3$ & branch on $\le 0$ & total given branches \\
\code{loop},\code{loop2} & $3,5$ & bounded iteration & terminating (resource: timeout) \\
\code{compr} & $2$ & unbounded search & \emph{genuine}: may not terminate \\
\code{push} & $2$ & stack push & total (resource: push limit) \\
\code{pop} & $1$ & stack pop & \emph{genuine}: error on empty \\
\bottomrule
\end{tabular}
\caption{The synthesis DSL. Only \code{divi}/\code{modu} (zero divisor), \code{pop}
(empty stack), and \code{compr} (non-terminating search) introduce \emph{genuine}
partiality; \code{loop}/\code{loop2}, overflow, and timeouts are \emph{resource} limits.
This three-way character drives the totality analysis.}
\label{tab:ops}
\end{table}

A closed form illustrates the language: the factorial $n!$ is
\code{loop(mult(x,y),\,x,\,one)} --- fold multiplication by the running counter $n$ times
starting from $1$. The counting sequence $a(n)=n+1$ is \code{addi(x,one)}.

\section{Verified GSLT semantics}
\label{sec:gslt}

The DSL is given a machine-checked operational semantics in Lean. On top of it we build a
generative-semantics layer as a triple $(T,E,R)$: terms $T$, an equational theory $E$
(a setoid), and a rewrite relation $R$ with left/right coherence laws. The coarse
\emph{big-step GSLT} rewrites a program to the integer value its fuel-bounded evaluator
emits; a completed-value term is never reflexively a program, so non-triviality is
load-bearing. Two modalities express synthesis goals:
\begin{align*}
  \Diamond(\text{emits } v \text{ at } k) &\iff \exists\, s',\ \code{eval}\ \mathit{fuel}\ \mathit{sig}\ p\ (\code{seed}\ k)\ = \code{some}(v,s'),\\
  \code{EmitsPrefix}\ p\ \ell &\iff \forall i\, v,\ \ell_i = v \Rightarrow \Diamond(\text{emits } v \text{ at } i).
\end{align*}
\code{EmitsPrefix} is precisely the OEIS synthesis specification: ``$p$ reproduces the
observed prefix $\ell$.''

\paragraph{Obligations.} A synthesizer that claims correctness must respect a family of
formal obligations. We single out three that are discharged as machine-checked theorems,
axiom-clean (each theorem's \code{\#print axioms} footprint lies within
$\{\text{propext}, \text{Quot.sound}\}$):
\begin{obligation}[Recognizer completeness / mask equivalence]
The legality predicate \code{maskLegal} on token streams coincides with the executable
recognizer (\code{recognize\_iff\_maskLegal}), and every well-formed program is recognized
(\code{recognize\_complete}). Hence a constrained decoder that emits only mask-legal tokens
can generate exactly the well-formed programs.
\end{obligation}
\begin{obligation}[Relabeling invariance]
Evaluation is invariant under a signature isomorphism
(\code{eval\_relabel\_signature\_iso}): renaming primitives consistently does not change the
computed sequence. Behaviour depends on structure, not on operator names.
\end{obligation}
\begin{obligation}[Canonicalization soundness]
A flag-driven canonical form (sorting the children of commutative operators, etc.) is
semantics-preserving (\code{canonicalize\_sound}, \code{orgCanonicalize\_sound}), so
deduplicating by canonical form never merges behaviourally distinct programs.
\end{obligation}
Further obligations (equational-theory soundness, eval-congruence for the modalities,
finite-probe rigidity, small-step/big-step agreement) organize the remaining work. Finite-
probe rigidity---whether agreement on a finite probe entails equality---is the formal
shadow of the leakage phenomenon in Section~\ref{sec:extensional}.

\section{The synthesizer as specification}
\label{sec:synth}

The neural synthesizer maps a sequence prefix to a program graph (sequence-to-graph). Its
decoder is \emph{execution-guided}: at each step a legality mask restricts the token set to
those that can extend a well-formed program, and the mask is exactly the recognizer of
Obligation~1. The training ladder A0--A3 rises from a plain conditioned decoder (A0) to
retrieval-augmented and contrastive variants (A3) that consult a similarity index. The
central design stance is that the network's operational contracts are the prover's
theorems: the mask's soundness, the canonical form's soundness, and relabeling invariance
are not engineering conveniences but the denotational safety envelope of the synthesizer.

\section{Leakage-safe data and the extensional layer}
\label{sec:extensional}

\paragraph{Exact deduplication and splits.} Each program's probe signature \sig{} is its
output vector on $n=0,\dots,31$ at a fixed fuel, with $\bot$ (\code{None}) marking a
timeout. Programs with identical \sig{} form an exact behavioural equivalence class
(``E-class''); the train/val/test split is assigned by whole connected component so that no
exact behavioural duplicate straddles the split boundary. Over $104{,}083$ programs this
gate is leak-tight for exact duplicates: all $3{,}007$ exact-\sig{} duplicate pairs lie in
the same split.

\paragraph{Graded similarity.} Exact equality is brittle. We express graded similarity as a
probabilistic-logic truth value $\langle s,c\rangle$~\cite{pln2009}: the extensional
strength is behavioural agreement over the \emph{defined} domain,
$s = |{\text{agree}}| / |{\text{informative}}|$, and confidence follows the evidence-count
law $c = m/(m+k)$ with $m$ the number of informative probe positions. Exact equality is the
degenerate case $s=1,\ m=32$. A locality-sensitive index over signatures generates
near-equivalent candidate pairs without an $O(n^2)$ scan.

\paragraph{Soft-leakage audit.} The graded layer surfaces \emph{near}-duplicates the exact
gate cannot see. Table~\ref{tab:leak} reports pairs whose two programs sit in different
splits. The benchmark-integrity figure is the fraction of evaluation programs with a near
behavioural twin in \emph{train}: at strength $\ge 0.95$, $1.84\%$ of test and $1.84\%$ of
val programs (lower bounds---the index can miss a pair, never invent one). A
disagreement-locus analysis separates genuine near-twins (disagreements confined to a
single term, none in the tail) from ``false friends'' that agree only on a shared trivial
tail; the trustworthy signal concentrates in the $\ge 0.95$ band. This vindicates gating on
\emph{exact} full-signature equality rather than a similarity threshold, which would either
over-merge on trivial tails or miss single-term twins.

\begin{table}[t]
\centering
\begin{tabular}{@{}lrr@{}}
\toprule
strength band & cross-split pairs & train$\leftrightarrow$test pairs \\
\midrule
$0.80$--$0.90$ & $17{,}250$ & $6{,}365$ \\
$0.90$--$0.95$ & $2{,}614$  & $909$ \\
$0.95$--$0.99$ & $502$      & $150$ \\
\bottomrule
\end{tabular}
\caption{Soft leakage: near-equivalent program pairs ($d\ge 1$, not byte-identical) placed
in different splits by the exact gate. Distinct evaluation programs with a train near-twin:
$1.84\%$ (test and val) at $\ge 0.95$; $5.4$--$5.8\%$ at $\ge 0.90$.}
\label{tab:leak}
\end{table}

\section{The intensional layer: provable properties}
\label{sec:intensional}

Where the extensional layer reads behaviour off finitely many outputs, the intensional
layer \emph{proves} properties of the function for all inputs by sound abstract
interpretation~\cite{cousot1977}. Each transfer function over-approximates the concrete
operator, so a claim of ``always $\ge 0$'' or ``always even'' holds by construction; loops
are handled by Kleene fixpoints in finite lattices. We compute three properties.

\begin{property}[Totality]
A program is \emph{provably total} when only resource limits can stop it: it contains
neither \code{compr} nor \code{pop}, and every \code{divi}/\code{modu} divisor is provably
nonzero. This is a conservative lower bound---``unknown'' does not mean partial.
\end{property}
\begin{property}[Sign]
A sign-domain analysis (with $\code{x}=n\ge 0$, and $a\bmod b\ge 0$ for $b>0$) proves the
OEIS keyword \code{nonn} (nonnegativity) for many programs.
\end{property}
\begin{property}[Parity]
A $\mathbb{Z}/2$ analysis proves ``always even'' / ``always odd'' (e.g.\ \code{mult} by an
even factor is even).
\end{property}

\paragraph{Soundness gates.} Two independent checks validate the implementation. A static
gate compares every sign/parity claim against the real signatures: $0$ violations across
all $104{,}083$ programs ($\approx 3.3$M value-checks). An empirical totality gate re-runs
the real evaluator on $500$ provably-total programs to $n=48$ (beyond the probe window):
$0$ genuine errors, only resource timeouts. Table~\ref{tab:props} gives the distribution.

\begin{table}[t]
\centering
\begin{tabular}{@{}lrr@{}}
\toprule
provable property & count & share \\
\midrule
total          & $37{,}996$ & $36.5\%$ \\
nonnegative (\code{nonn}) & $45{,}251$ & $43.5\%$ \\
positive       & $14{,}907$ & $14.3\%$ \\
even           & $2{,}725$  & $2.6\%$ \\
odd            & $2{,}926$  & $2.8\%$ \\
\bottomrule
\end{tabular}
\caption{Provable-property distribution over $104{,}083$ programs. Shares are lower bounds:
the analysis is conservative.}
\label{tab:props}
\end{table}

\paragraph{Disambiguating the signature.} A $\bot$ in \sig{} is ambiguous: is the program
genuinely partial, or merely slow at the probe budget? Totality resolves it. Of the
$12{,}966$ programs carrying a $\bot$, $2{,}345$ are provably total---their $\bot$ is a
\emph{resource artefact}, not partiality---while $10{,}621$ are not provably total, with
attributed cause: $6{,}539$ \code{compr} (non-terminating search), $3{,}023$ \code{pop}
(empty list), $1{,}059$ \code{divi}/\code{modu} (possible zero divisor). The intensional
layer thereby cleans the extensional signal.

\section{Filtering and measuring in the OEIS loop}
\label{sec:filter}

Both roles Table~\ref{tab:props} enables are useful inside a synthesis loop.

\paragraph{Filtering.} An OEIS target often carries known properties---its keywords
(\code{nonn}, \code{sign}, \code{mult}), or properties provable from a few terms. A
candidate program whose \emph{provable} properties contradict the target's is provably
wrong and can be pruned before it is ever run: if the target is nonnegative, every
provably-negative candidate is rejected soundly. Lifted into the decoder, this is a
\emph{verified semantic mask} on top of the syntactic legality mask of
Section~\ref{sec:synth}: constrain generation to programs provably producing outputs of the
required sign/parity. Because the properties are sound, such filtering never discards a
correct program.

\paragraph{Measuring.} The same properties are theoretical measurements of a program and of
a sequence. Totality distinguishes a genuinely partial sequence from a slow-to-compute total
one; the sign/parity certificates are machine-checked versions of OEIS keywords, so the
corpus can be scored---and OEIS annotations predicted and \emph{verified}---at scale. More
finely, the graded truth values give each program pair a calibrated similarity, and the
provable properties give certain-confidence intensional features that pool with behavioural
evidence by probabilistic-logic revision. The corpus becomes not merely programs and
outputs but programs, outputs, and \emph{proofs}.

\section{The loop compounds: a three-seed cumulative replication}
\label{sec:loop}

The construction above is only interesting if the loop it serves actually runs. It does. We
report a prospectively registered replication of the self-learning regime of
\cite{gauthier2023} using a typed-action decoder over the verified skeleton of
Section~\ref{sec:synth}: the network never emits program text, only \emph{legal construction
actions} over an exact state, and the checker disposes.

The protocol is the cumulative one. Each of three independently seeded worlds---seeds are
random-number labels only, controlling initialisation, minibatch order, dropout, and search
tie-breaking; the data pool ($79{,}249$ rows) and the $240$ search tasks are shared---trains
on its accumulated verified corpus, searches under a fixed budget, checks every candidate
against the whole OEIS, adds newly covered sequences to its corpus, and repeats. The
endpoint is cumulative coverage $C_g$ and its increment $D_g = C_g - C_{g-1}$; no
comparison against held-back external data governs continuation, since in the deployed
regime no such reservoir exists.

The shared initial corpus is not chosen freely: it is the reference system's own public
bootstrap seed \code{sol0}~\cite{gauthier2023}, $3{,}771$ sequences found by \emph{random}
program search with no learning, and the very origin from which that system subsequently
grew its full published corpus through many iterations of the loop we are running here.
Starting from that seed rather than from a mature corpus is deliberate, with three
consequences. The origin is reproducible and learning-free, a fixed public artefact anyone
can replay from. Our run is then the \emph{same} self-learning process re-executed---not a
fork of another system's answers---so a rediscovery is a genuine held-out generalisation
within our own lineage rather than a lookup. And the base stays small enough that a world's
own discoveries form a measurable share of the training signal instead of being drowned by
it, which is what makes the loop's dynamics observable at all. Coverage per world:

\begin{center}
\begin{tabular}{lrrr}
\toprule
world & initial & after gen.\ 1 & after gen.\ 2 \\
\midrule
17 & 3{,}771 & 3{,}789 & 3{,}846 \\
29 & 3{,}771 & 3{,}841 & 3{,}891 \\
43 & 3{,}771 & 3{,}807 & 3{,}862 \\
\bottomrule
\end{tabular}
\end{center}

At generation three we registered, before any update, a three-way comparison of how the
accumulated discoveries enter training: \emph{uniform} replay (literal pooling, under which
the $194$ cumulative discoveries receive their natural $0.24\%$ of sampling mass);
\emph{balanced} replay ($25\%$ discovery mass); and balanced replay with the auxiliary
predictive heads disabled. All three arms continue from the same per-world parent
checkpoint under identical update counts and search budgets, and are judged solely by new
whole-OEIS A-numbers beyond the pooled $3{,}986$ already known. New A-numbers, generation
three:

\begin{center}
\begin{tabular}{lrrrr}
\toprule
setting & world 17 & world 29 & world 43 & pooled \\
\midrule
uniform replay & 48 & 39 & 72 & 159 \\
balanced replay, aux.\ off & 65 & 155 & 88 & 308 \\
balanced replay, aux.\ on & 135 & 155 & 177 & \textbf{467} \\
\bottomrule
\end{tabular}
\end{center}

Every setting is positive in every world: nine of nine cells extend coverage, so the loop
compounds rather than saturating at this scale. Two registered contrasts separate the
mechanism. Balanced replay is worth $+308$ pooled over uniform---at literal pooling the
model's own discoveries are numerically drowned by the base corpus, and the loop is starved
of exactly the signal it exists to generate. The auxiliary heads are worth a further $+159$.
The generation-three increment under the selected setting exceeds the generation-two
increment by roughly threefold, and the setting carried into generation four was fixed by
the registered selection rule rather than chosen after inspection.

Two negative results bound the claim, and we record them because they cost us the
distinction. First, an ablation asking whether a world's own discoveries beat an equal mass
of \emph{unused external certified examples} answers no, in all three worlds
($162$ versus $253$ pooled): feedback submitted more candidates yet explored fewer distinct
programs ($5{,}540$ versus $6{,}233$), a diversity contraction rather than a fitting
failure. That comparison is a legitimate ablation and a poor objective---the fresh arm is an
oracle-advantaged comparator that does not exist in the deployed loop, where the corpus is
whatever the loop has verified. Second, no arm in any generation solved its own assigned
held-out target within budget. The evidence therefore concerns learned \emph{discovery}
guidance over the catalogue---the regime \cite{gauthier2023} operates in---and not
target-conditioned synthesis, which remains open.

What the verified layer contributes here is attributability rather than yield. Because the
legality mask is machine-checked, a discovery cannot be an artefact of a malformed
candidate slipping past a hand-written filter; because ranking is provably
recall-preserving, guidance changes search order and cannot manufacture or destroy
solutions; and because every accepted program is checked against the whole corpus, the
increments above are exact rather than estimated. The loop's growth and the loop's
trustworthiness are established by different machinery, and both are needed.

\section{From leakage to conjecture}
\label{sec:conjecture}

Near-equivalence is dual to conjecture: what threatens a benchmark as contamination is, read
the other way, a discovered near-identity. The sharpest instance is the pairing of the
\emph{fastest} and the \emph{smallest} program for a sequence. These minimize opposite terms
of Levin's $Kt$ complexity, $Kt(x)=\min_p |p|+\log t(p)$~\cite{levin1973}: program length
and running time. Two quasi-independent searches, under different inductive biases,
converging on the same total function is strong corroboration that the function is the
intended sequence---never entailment (finitely many terms cannot pin an infinite sequence),
but Bayesian evidence, and in probabilistic-logic terms the revision of two independent
estimates into higher confidence. \emph{Proving} the two programs equivalent on all $n$
converts finite-prefix agreement into total agreement; the induction predicates such proofs
require are exactly what recent conjecturing systems learn to
generate~\cite{gauthier2025}. Our verified layer offers to discharge them as kernel-checked
theorems rather than trusted external calls. The interesting conjectures live where
extensional similarity is high but \emph{syntactic} intensional similarity is low---two
unlike programs computing one sequence---and shared \emph{provable} properties become lemmas
toward the identity proof.

\section{Status and roadmap}
\label{sec:roadmap}

The verified DSL semantics, recognizer, big-step modalities, and the three crown obligations
are machine-checked and axiom-clean. The synthesizer's data pipeline, exact split gate, and
the WM-PLN extensional and intensional layers are implemented and measured over the current
corpus, with the numbers reported above; the self-learning loop of Section~\ref{sec:loop}
compounds under three registered replay settings across three worlds. Near-term work: (i) bind the Lean canonical policy
to the exported operator-table metadata so the ML mask and the proved canonical form share
one source of truth; (ii) lift the sound Python properties to Lean theorems, closing the
verified-semantic-mask loop; (iii) retain the fastest and smallest program per sequence so
the within-sequence equivalence-conjecture analysis of Section~\ref{sec:conjecture} becomes
runnable; (iv) add growth-class and monotonicity to the abstract interpreter. The unifying
aim is an OEIS synthesis loop in which every accepted program is both behaviourally
corroborated and provably constrained.

\begin{thebibliography}{9}
\bibitem{gauthier2023} T.~Gauthier and J.~Urban.
  \emph{Learning Program Synthesis for Integer Sequences from Scratch}.
  AAAI 2023; arXiv:2202.11908.
\bibitem{gauthier2025} T.~Gauthier and J.~Urban.
  \emph{Learning Conjecturing from Scratch}. arXiv:2503.01389, 2025.
\bibitem{oeis} OEIS Foundation Inc.
  \emph{The On-Line Encyclopedia of Integer Sequences}. \url{https://oeis.org}.
\bibitem{cousot1977} P.~Cousot and R.~Cousot.
  \emph{Abstract Interpretation: A Unified Lattice Model for Static Analysis of Programs by
  Construction or Approximation of Fixpoints}. POPL 1977.
\bibitem{levin1973} L.~A.~Levin.
  \emph{Universal Sequential Search Problems}. Problems of Information Transmission,
  9(3):265--266, 1973.
\bibitem{pln2009} B.~Goertzel, M.~Ikl\'e, I.~Goertzel, and A.~Heljakka.
  \emph{Probabilistic Logic Networks}. Springer, 2009.
\end{thebibliography}

\end{document}
